This work restructures the original CPE code release \citep{mack2026} into a library.
Two changes matter for the results that follow: the search over the factor population is cheaper, and the implementation reproduces the paper.
The engineering that is not specific to the method, packaging it for import and running experiments on rented hardware, is described in Appendices~\ref{app:package} and~\ref{app:remote}.

\subsection{Staggered generation search}

The generation process now does a more efficient staggered search over the factor population, successively halving the search space, instead of always generating the full sweep.
Each round scores the surviving factors on the same small set of prompts, prunes the weakest, and gives the survivors more prompts.
This is not always possible.
When the property of interest is a property of the population, as when investigating conversational coherence, the full sweep is still needed.
When looking for one specific behaviour, as in the sandbagging case, the staggered search finds it at a fraction of the generation cost.

\subsection{Replication of the paper's results}

We have ensured that the new implementation works as well as the old one by replicating a few key results of the paper.
Table~\ref{tab:replication} compares our scores against the published ones on the paper's natural-behaviour environments, together with its two unsupervised controls, random LoRAs at the same locations and SAE steering.

\begin{table}[htbp]
  \centering
  \small
  % Generated by exhibits/replication/table_replication.py -- do not edit by hand.
\begin{tabular}{llrrrr}
\toprule
Environment & Model & Base & Random-LoRA & SAE & CPE \\
\midrule
Countdown & Qwen3-8B & 0.77 (0.73) & 0.81 (0.75) & --- (0.78) & \textbf{0.88 (0.85)} \\
Countdown & Llama-3.1-8B & 0.15 (0.07) & 0.11 (0.08) & 0.32 (0.32) & \textbf{0.51 (0.36)} \\
Sycophancy & Qwen3-8B & 0.88 (0.88) & 0.88 (0.88) & 0.96 (0.92) & \textbf{0.90 (0.96)} \\
Sycophancy & Llama-3.1-8B & 0.80 (0.79) & 0.80 (0.79) & 0.80 (0.82) & \textbf{0.92 (0.93)} \\
Jailbreak (ASR) & robust-llama3-8b & 0.00 (0.00) & 0.00 (0.00) & 0.00 (0.00) & \textbf{0.52 (0.65)} \\
\bottomrule
\end{tabular}

  \caption{%
    Replication of the CPE paper's Table~1.
    Each cell gives our score with the published value in parentheses.
    Scores are the selected top-1 factor's held-out test score; base is the same run's no-adapter baseline.
    Countdown and sycophancy are accuracies, jailbreak is attack success rate.%
  }
  \label{tab:replication}
\end{table}

CPE reproduces on every environment.
The random-LoRA control sits at baseline throughout, confirming that the gains come from the trained direction rather than from perturbing weights at all.
Effect sizes are in line with or exceed the published ones, with the jailbreak environment the one place we come in under, while still reproducing its qualitative headline: an attack success rate of zero rises to about half against a model adversarially trained to resist latent attacks.
